\documentclass[aspectratio=169, 11pt]{beamer}
% ============================================================
% Preambles/header.tex
% MN52233 — Causal Identification Strategies
% University of Bath, School of Management
% ============================================================
% Shared preamble for all lecture slide decks.
% Usage in each .tex file: % ============================================================
% Preambles/header.tex
% MN52233 — Causal Identification Strategies
% University of Bath, School of Management
% ============================================================
% Shared preamble for all lecture slide decks.
% Usage in each .tex file: % ============================================================
% Preambles/header.tex
% MN52233 — Causal Identification Strategies
% University of Bath, School of Management
% ============================================================
% Shared preamble for all lecture slide decks.
% Usage in each .tex file: \input{../Preambles/header.tex}
% ============================================================

% --- Core packages ---
\usepackage{amsmath,amssymb,mathtools}
\usepackage{xcolor}
\usepackage{graphicx}
\usepackage{booktabs}
\usepackage{tikz}
\usetikzlibrary{arrows.meta,positioning,shapes.geometric,calc,decorations.pathreplacing}
\usepackage{tcolorbox}
\tcbuselibrary{skins}
\usepackage{natbib}
\bibpunct{(}{)}{;}{a}{,}{,}
\usepackage{enumerate}
\usepackage{ragged2e}
\usepackage{microtype}

% ============================================================
% COLOUR DEFINITIONS
% ============================================================

% Primary palette
\definecolor{bathnav}{HTML}{1F3A5F}      % Dark navy — primary
\definecolor{bathacc}{HTML}{3A7EBF}      % Medium blue — secondary
\definecolor{bathlight}{HTML}{EDF3FB}    % Very light blue — box backgrounds
\definecolor{nearblack}{HTML}{1A1A1A}    % Near-black body text

% Okabe-Ito colorblind-safe palette (mandatory for all figures)
\definecolor{okabeblue}{HTML}{0072B2}
\definecolor{okabeorange}{HTML}{E69F00}
\definecolor{okabegreen}{HTML}{009E73}
\definecolor{okabesky}{HTML}{56B4E9}
\definecolor{okabevermillion}{HTML}{D55E00}
\definecolor{okabepink}{HTML}{CC79A7}
\definecolor{okabeyellow}{HTML}{F0E442}

% ============================================================
% BEAMER THEME CUSTOMISATION
% ============================================================

\usetheme{default}

% --- Text & structure colours ---
\setbeamercolor{normal text}{fg=nearblack, bg=white}
\setbeamercolor{structure}{fg=bathnav}
\setbeamercolor{alerted text}{fg=okabevermillion}
\setbeamercolor{example text}{fg=okabegreen}

% --- Title page ---
\setbeamercolor{title}{fg=bathnav}
\setbeamercolor{subtitle}{fg=bathacc}
\setbeamercolor{author}{fg=nearblack}
\setbeamercolor{institute}{fg=nearblack!70}
\setbeamercolor{date}{fg=bathacc}

% --- Frame title ---
\setbeamercolor{frametitle}{bg=bathnav, fg=white}
\setbeamerfont{frametitle}{series=\bfseries, size=\normalsize}

% --- Itemize bullets ---
\setbeamercolor{item}{fg=bathnav}
\setbeamercolor{subitem}{fg=bathacc}
\setbeamertemplate{itemize item}{\small\raise0.5pt\hbox{\color{bathnav}$\bullet$}}
\setbeamertemplate{itemize subitem}{\small\raise0.5pt\hbox{\color{bathacc}$\circ$}}
\setbeamertemplate{itemize subsubitem}{\scriptsize\raise1pt\hbox{\color{bathacc}$-$}}

% --- Block colours ---
\setbeamercolor{block title}{bg=bathnav, fg=white}
\setbeamercolor{block body}{bg=bathlight, fg=nearblack}
\setbeamercolor{block title alerted}{bg=okabevermillion, fg=white}
\setbeamercolor{block body alerted}{bg=okabevermillion!10, fg=nearblack}
\setbeamercolor{block title example}{bg=okabegreen, fg=white}
\setbeamercolor{block body example}{bg=okabegreen!10, fg=nearblack}

% --- Frame title template ---
\setbeamertemplate{frametitle}{%
  \nointerlineskip%
  \begin{beamercolorbox}[wd=\paperwidth, ht=2.8ex, dp=1.2ex, leftskip=1em]{frametitle}%
    \usebeamerfont{frametitle}\insertframetitle%
  \end{beamercolorbox}%
  \nointerlineskip%
  {\color{bathacc}\rule{\paperwidth}{0.4pt}}%
  \vspace{0.3em}%
}

% --- Remove navigation symbols ---
\setbeamertemplate{navigation symbols}{}

% --- Footer ---
\setbeamercolor{footline}{fg=nearblack!60, bg=white}
\setbeamertemplate{footline}{%
  \leavevmode%
  {\color{bathnav!30}\rule{\paperwidth}{0.3pt}}%
  \hbox{%
    \begin{beamercolorbox}[wd=.30\paperwidth, ht=2.2ex, dp=1ex, leftskip=0.8em]{footline}%
      \footnotesize\color{bathacc}MN52233%
    \end{beamercolorbox}%
    \begin{beamercolorbox}[wd=.40\paperwidth, ht=2.2ex, dp=1ex, center]{footline}%
      \footnotesize\color{nearblack!60}\insertshorttitle%
    \end{beamercolorbox}%
    \begin{beamercolorbox}[wd=.30\paperwidth, ht=2.2ex, dp=1ex, rightskip=0.8em, right]{footline}%
      \footnotesize\color{nearblack!60}\insertframenumber\,/\,\inserttotalframenumber%
    \end{beamercolorbox}%
  }%
  \vskip2pt%
}

% --- Title page template ---
\setbeamertemplate{title page}{%
  \vspace{1.8em}%
  \begin{beamercolorbox}[sep=6pt, center]{title}%
    \usebeamerfont{title}\usebeamercolor[fg]{title}\inserttitle\par%
  \end{beamercolorbox}%
  \vskip0.4em%
  \begin{beamercolorbox}[sep=4pt, center]{subtitle}%
    \usebeamerfont{subtitle}\usebeamercolor[fg]{subtitle}\insertsubtitle\par%
  \end{beamercolorbox}%
  \vskip0.8em%
  {\color{bathnav!40}\rule{0.5\paperwidth}{0.4pt}}\par%
  \vskip0.8em%
  \begin{beamercolorbox}[sep=3pt, center]{author}%
    \usebeamerfont{author}\insertauthor%
  \end{beamercolorbox}%
  \begin{beamercolorbox}[sep=2pt, center]{institute}%
    \usebeamerfont{institute}\insertinstitute%
  \end{beamercolorbox}%
  \vskip0.4em%
  \begin{beamercolorbox}[sep=2pt, center]{date}%
    \usebeamerfont{date}\insertdate%
  \end{beamercolorbox}%
}

% ============================================================
% SECTION TRANSITION SLIDES
% ============================================================

% Usage: \sectionslide{Section Title}{Optional subtitle or question}
\newcommand{\sectionslide}[2]{%
  \begingroup%
  \setbeamercolor{background canvas}{bg=bathnav}%
  \begin{frame}[plain, noframenumbering]%
    \vfill%
    \centering%
    {\Large\bfseries\textcolor{white}{#1}}\par%
    \vskip0.8em%
    \ifx\relax#2\relax\else%
      {\normalsize\textcolor{okabesky}{#2}}\par%
    \fi%
    \vfill%
  \end{frame}%
  \endgroup%
}

% ============================================================
% TCOLORBOX CUSTOM ENVIRONMENTS
% ============================================================

\tcbset{
  beamerbase/.style={
    enhanced,
    boxrule=0.4pt,
    arc=2pt,
    left=6pt, right=6pt, top=3pt, bottom=4pt,
    fonttitle=\bfseries\small,
    attach boxed title to top left={yshift=-0.1mm},
    boxed title style={
      arc=2pt,
      boxrule=0pt,
      left=6pt, right=6pt, top=3pt, bottom=3pt,
    },
    before upper={\setlength\parindent{0pt}\small},
  }
}

% Key point box — dark navy title bar, white title text, light blue body
\newtcolorbox{keybox}[1][Key Point]{
  beamerbase,
  colback=bathnav!8,
  colframe=bathnav,
  colbacktitle=bathnav,
  coltitle=white,
  title={#1},
}

% Paper example box — dark orange title bar, white title text, light orange body
\newtcolorbox{examplebox}[1][Example]{
  beamerbase,
  colback=okabeorange!12,
  colframe=okabeorange!80!black,
  colbacktitle=okabeorange!85!black,
  coltitle=white,
  title={#1},
}

% Warning / caveat box — vermillion title bar, white title text, light red body
\newtcolorbox{warningbox}[1][Caution]{
  beamerbase,
  colback=okabevermillion!10,
  colframe=okabevermillion,
  colbacktitle=okabevermillion,
  coltitle=white,
  title={#1},
}

% Intuition box — green title bar, white title text, light green body
\newtcolorbox{intuitionbox}[1][Intuition]{
  beamerbase,
  colback=okabegreen!10,
  colframe=okabegreen,
  colbacktitle=okabegreen,
  coltitle=white,
  title={#1},
}

% ============================================================
% MATH MACROS
% ============================================================

\newcommand{\E}[1]{\mathbb{E}\!\left[#1\right]}
\newcommand{\Econd}[2]{\mathbb{E}\!\left[#1\,\middle|\,#2\right]}
\newcommand{\Yi}[1]{Y_{i}(#1)}
\newcommand{\Y}[1]{Y(#1)}
\newcommand{\ATE}{\mathrm{ATE}}
\newcommand{\ATT}{\mathrm{ATT}}
\newcommand{\ITT}{\mathrm{ITT}}
\newcommand{\LATE}{\mathrm{LATE}}
\newcommand{\indep}{\perp\!\!\!\perp}
\newcommand{\muted}[1]{{\color{nearblack!50}#1}}
\newcommand{\highlight}[1]{{\color{bathnav}\textbf{#1}}}
\newcommand{\source}[1]{\vspace{0.3em}\par{\scriptsize\color{nearblack!60}\textit{#1}}}

% ============================================================
% BIBLIOGRAPHY STYLE
% ============================================================

\bibliographystyle{plainnat}

% ============================================================

% --- Core packages ---
\usepackage{amsmath,amssymb,mathtools}
\usepackage{xcolor}
\usepackage{graphicx}
\usepackage{booktabs}
\usepackage{tikz}
\usetikzlibrary{arrows.meta,positioning,shapes.geometric,calc,decorations.pathreplacing}
\usepackage{tcolorbox}
\tcbuselibrary{skins}
\usepackage{natbib}
\bibpunct{(}{)}{;}{a}{,}{,}
\usepackage{enumerate}
\usepackage{ragged2e}
\usepackage{microtype}

% ============================================================
% COLOUR DEFINITIONS
% ============================================================

% Primary palette
\definecolor{bathnav}{HTML}{1F3A5F}      % Dark navy — primary
\definecolor{bathacc}{HTML}{3A7EBF}      % Medium blue — secondary
\definecolor{bathlight}{HTML}{EDF3FB}    % Very light blue — box backgrounds
\definecolor{nearblack}{HTML}{1A1A1A}    % Near-black body text

% Okabe-Ito colorblind-safe palette (mandatory for all figures)
\definecolor{okabeblue}{HTML}{0072B2}
\definecolor{okabeorange}{HTML}{E69F00}
\definecolor{okabegreen}{HTML}{009E73}
\definecolor{okabesky}{HTML}{56B4E9}
\definecolor{okabevermillion}{HTML}{D55E00}
\definecolor{okabepink}{HTML}{CC79A7}
\definecolor{okabeyellow}{HTML}{F0E442}

% ============================================================
% BEAMER THEME CUSTOMISATION
% ============================================================

\usetheme{default}

% --- Text & structure colours ---
\setbeamercolor{normal text}{fg=nearblack, bg=white}
\setbeamercolor{structure}{fg=bathnav}
\setbeamercolor{alerted text}{fg=okabevermillion}
\setbeamercolor{example text}{fg=okabegreen}

% --- Title page ---
\setbeamercolor{title}{fg=bathnav}
\setbeamercolor{subtitle}{fg=bathacc}
\setbeamercolor{author}{fg=nearblack}
\setbeamercolor{institute}{fg=nearblack!70}
\setbeamercolor{date}{fg=bathacc}

% --- Frame title ---
\setbeamercolor{frametitle}{bg=bathnav, fg=white}
\setbeamerfont{frametitle}{series=\bfseries, size=\normalsize}

% --- Itemize bullets ---
\setbeamercolor{item}{fg=bathnav}
\setbeamercolor{subitem}{fg=bathacc}
\setbeamertemplate{itemize item}{\small\raise0.5pt\hbox{\color{bathnav}$\bullet$}}
\setbeamertemplate{itemize subitem}{\small\raise0.5pt\hbox{\color{bathacc}$\circ$}}
\setbeamertemplate{itemize subsubitem}{\scriptsize\raise1pt\hbox{\color{bathacc}$-$}}

% --- Block colours ---
\setbeamercolor{block title}{bg=bathnav, fg=white}
\setbeamercolor{block body}{bg=bathlight, fg=nearblack}
\setbeamercolor{block title alerted}{bg=okabevermillion, fg=white}
\setbeamercolor{block body alerted}{bg=okabevermillion!10, fg=nearblack}
\setbeamercolor{block title example}{bg=okabegreen, fg=white}
\setbeamercolor{block body example}{bg=okabegreen!10, fg=nearblack}

% --- Frame title template ---
\setbeamertemplate{frametitle}{%
  \nointerlineskip%
  \begin{beamercolorbox}[wd=\paperwidth, ht=2.8ex, dp=1.2ex, leftskip=1em]{frametitle}%
    \usebeamerfont{frametitle}\insertframetitle%
  \end{beamercolorbox}%
  \nointerlineskip%
  {\color{bathacc}\rule{\paperwidth}{0.4pt}}%
  \vspace{0.3em}%
}

% --- Remove navigation symbols ---
\setbeamertemplate{navigation symbols}{}

% --- Footer ---
\setbeamercolor{footline}{fg=nearblack!60, bg=white}
\setbeamertemplate{footline}{%
  \leavevmode%
  {\color{bathnav!30}\rule{\paperwidth}{0.3pt}}%
  \hbox{%
    \begin{beamercolorbox}[wd=.30\paperwidth, ht=2.2ex, dp=1ex, leftskip=0.8em]{footline}%
      \footnotesize\color{bathacc}MN52233%
    \end{beamercolorbox}%
    \begin{beamercolorbox}[wd=.40\paperwidth, ht=2.2ex, dp=1ex, center]{footline}%
      \footnotesize\color{nearblack!60}\insertshorttitle%
    \end{beamercolorbox}%
    \begin{beamercolorbox}[wd=.30\paperwidth, ht=2.2ex, dp=1ex, rightskip=0.8em, right]{footline}%
      \footnotesize\color{nearblack!60}\insertframenumber\,/\,\inserttotalframenumber%
    \end{beamercolorbox}%
  }%
  \vskip2pt%
}

% --- Title page template ---
\setbeamertemplate{title page}{%
  \vspace{1.8em}%
  \begin{beamercolorbox}[sep=6pt, center]{title}%
    \usebeamerfont{title}\usebeamercolor[fg]{title}\inserttitle\par%
  \end{beamercolorbox}%
  \vskip0.4em%
  \begin{beamercolorbox}[sep=4pt, center]{subtitle}%
    \usebeamerfont{subtitle}\usebeamercolor[fg]{subtitle}\insertsubtitle\par%
  \end{beamercolorbox}%
  \vskip0.8em%
  {\color{bathnav!40}\rule{0.5\paperwidth}{0.4pt}}\par%
  \vskip0.8em%
  \begin{beamercolorbox}[sep=3pt, center]{author}%
    \usebeamerfont{author}\insertauthor%
  \end{beamercolorbox}%
  \begin{beamercolorbox}[sep=2pt, center]{institute}%
    \usebeamerfont{institute}\insertinstitute%
  \end{beamercolorbox}%
  \vskip0.4em%
  \begin{beamercolorbox}[sep=2pt, center]{date}%
    \usebeamerfont{date}\insertdate%
  \end{beamercolorbox}%
}

% ============================================================
% SECTION TRANSITION SLIDES
% ============================================================

% Usage: \sectionslide{Section Title}{Optional subtitle or question}
\newcommand{\sectionslide}[2]{%
  \begingroup%
  \setbeamercolor{background canvas}{bg=bathnav}%
  \begin{frame}[plain, noframenumbering]%
    \vfill%
    \centering%
    {\Large\bfseries\textcolor{white}{#1}}\par%
    \vskip0.8em%
    \ifx\relax#2\relax\else%
      {\normalsize\textcolor{okabesky}{#2}}\par%
    \fi%
    \vfill%
  \end{frame}%
  \endgroup%
}

% ============================================================
% TCOLORBOX CUSTOM ENVIRONMENTS
% ============================================================

\tcbset{
  beamerbase/.style={
    enhanced,
    boxrule=0.4pt,
    arc=2pt,
    left=6pt, right=6pt, top=3pt, bottom=4pt,
    fonttitle=\bfseries\small,
    attach boxed title to top left={yshift=-0.1mm},
    boxed title style={
      arc=2pt,
      boxrule=0pt,
      left=6pt, right=6pt, top=3pt, bottom=3pt,
    },
    before upper={\setlength\parindent{0pt}\small},
  }
}

% Key point box — dark navy title bar, white title text, light blue body
\newtcolorbox{keybox}[1][Key Point]{
  beamerbase,
  colback=bathnav!8,
  colframe=bathnav,
  colbacktitle=bathnav,
  coltitle=white,
  title={#1},
}

% Paper example box — dark orange title bar, white title text, light orange body
\newtcolorbox{examplebox}[1][Example]{
  beamerbase,
  colback=okabeorange!12,
  colframe=okabeorange!80!black,
  colbacktitle=okabeorange!85!black,
  coltitle=white,
  title={#1},
}

% Warning / caveat box — vermillion title bar, white title text, light red body
\newtcolorbox{warningbox}[1][Caution]{
  beamerbase,
  colback=okabevermillion!10,
  colframe=okabevermillion,
  colbacktitle=okabevermillion,
  coltitle=white,
  title={#1},
}

% Intuition box — green title bar, white title text, light green body
\newtcolorbox{intuitionbox}[1][Intuition]{
  beamerbase,
  colback=okabegreen!10,
  colframe=okabegreen,
  colbacktitle=okabegreen,
  coltitle=white,
  title={#1},
}

% ============================================================
% MATH MACROS
% ============================================================

\newcommand{\E}[1]{\mathbb{E}\!\left[#1\right]}
\newcommand{\Econd}[2]{\mathbb{E}\!\left[#1\,\middle|\,#2\right]}
\newcommand{\Yi}[1]{Y_{i}(#1)}
\newcommand{\Y}[1]{Y(#1)}
\newcommand{\ATE}{\mathrm{ATE}}
\newcommand{\ATT}{\mathrm{ATT}}
\newcommand{\ITT}{\mathrm{ITT}}
\newcommand{\LATE}{\mathrm{LATE}}
\newcommand{\indep}{\perp\!\!\!\perp}
\newcommand{\muted}[1]{{\color{nearblack!50}#1}}
\newcommand{\highlight}[1]{{\color{bathnav}\textbf{#1}}}
\newcommand{\source}[1]{\vspace{0.3em}\par{\scriptsize\color{nearblack!60}\textit{#1}}}

% ============================================================
% BIBLIOGRAPHY STYLE
% ============================================================

\bibliographystyle{plainnat}

% ============================================================

% --- Core packages ---
\usepackage{amsmath,amssymb,mathtools}
\usepackage{xcolor}
\usepackage{graphicx}
\usepackage{booktabs}
\usepackage{tikz}
\usetikzlibrary{arrows.meta,positioning,shapes.geometric,calc,decorations.pathreplacing}
\usepackage{tcolorbox}
\tcbuselibrary{skins}
\usepackage{natbib}
\bibpunct{(}{)}{;}{a}{,}{,}
\usepackage{enumerate}
\usepackage{ragged2e}
\usepackage{microtype}

% ============================================================
% COLOUR DEFINITIONS
% ============================================================

% Primary palette
\definecolor{bathnav}{HTML}{1F3A5F}      % Dark navy — primary
\definecolor{bathacc}{HTML}{3A7EBF}      % Medium blue — secondary
\definecolor{bathlight}{HTML}{EDF3FB}    % Very light blue — box backgrounds
\definecolor{nearblack}{HTML}{1A1A1A}    % Near-black body text

% Okabe-Ito colorblind-safe palette (mandatory for all figures)
\definecolor{okabeblue}{HTML}{0072B2}
\definecolor{okabeorange}{HTML}{E69F00}
\definecolor{okabegreen}{HTML}{009E73}
\definecolor{okabesky}{HTML}{56B4E9}
\definecolor{okabevermillion}{HTML}{D55E00}
\definecolor{okabepink}{HTML}{CC79A7}
\definecolor{okabeyellow}{HTML}{F0E442}

% ============================================================
% BEAMER THEME CUSTOMISATION
% ============================================================

\usetheme{default}

% --- Text & structure colours ---
\setbeamercolor{normal text}{fg=nearblack, bg=white}
\setbeamercolor{structure}{fg=bathnav}
\setbeamercolor{alerted text}{fg=okabevermillion}
\setbeamercolor{example text}{fg=okabegreen}

% --- Title page ---
\setbeamercolor{title}{fg=bathnav}
\setbeamercolor{subtitle}{fg=bathacc}
\setbeamercolor{author}{fg=nearblack}
\setbeamercolor{institute}{fg=nearblack!70}
\setbeamercolor{date}{fg=bathacc}

% --- Frame title ---
\setbeamercolor{frametitle}{bg=bathnav, fg=white}
\setbeamerfont{frametitle}{series=\bfseries, size=\normalsize}

% --- Itemize bullets ---
\setbeamercolor{item}{fg=bathnav}
\setbeamercolor{subitem}{fg=bathacc}
\setbeamertemplate{itemize item}{\small\raise0.5pt\hbox{\color{bathnav}$\bullet$}}
\setbeamertemplate{itemize subitem}{\small\raise0.5pt\hbox{\color{bathacc}$\circ$}}
\setbeamertemplate{itemize subsubitem}{\scriptsize\raise1pt\hbox{\color{bathacc}$-$}}

% --- Block colours ---
\setbeamercolor{block title}{bg=bathnav, fg=white}
\setbeamercolor{block body}{bg=bathlight, fg=nearblack}
\setbeamercolor{block title alerted}{bg=okabevermillion, fg=white}
\setbeamercolor{block body alerted}{bg=okabevermillion!10, fg=nearblack}
\setbeamercolor{block title example}{bg=okabegreen, fg=white}
\setbeamercolor{block body example}{bg=okabegreen!10, fg=nearblack}

% --- Frame title template ---
\setbeamertemplate{frametitle}{%
  \nointerlineskip%
  \begin{beamercolorbox}[wd=\paperwidth, ht=2.8ex, dp=1.2ex, leftskip=1em]{frametitle}%
    \usebeamerfont{frametitle}\insertframetitle%
  \end{beamercolorbox}%
  \nointerlineskip%
  {\color{bathacc}\rule{\paperwidth}{0.4pt}}%
  \vspace{0.3em}%
}

% --- Remove navigation symbols ---
\setbeamertemplate{navigation symbols}{}

% --- Footer ---
\setbeamercolor{footline}{fg=nearblack!60, bg=white}
\setbeamertemplate{footline}{%
  \leavevmode%
  {\color{bathnav!30}\rule{\paperwidth}{0.3pt}}%
  \hbox{%
    \begin{beamercolorbox}[wd=.30\paperwidth, ht=2.2ex, dp=1ex, leftskip=0.8em]{footline}%
      \footnotesize\color{bathacc}MN52233%
    \end{beamercolorbox}%
    \begin{beamercolorbox}[wd=.40\paperwidth, ht=2.2ex, dp=1ex, center]{footline}%
      \footnotesize\color{nearblack!60}\insertshorttitle%
    \end{beamercolorbox}%
    \begin{beamercolorbox}[wd=.30\paperwidth, ht=2.2ex, dp=1ex, rightskip=0.8em, right]{footline}%
      \footnotesize\color{nearblack!60}\insertframenumber\,/\,\inserttotalframenumber%
    \end{beamercolorbox}%
  }%
  \vskip2pt%
}

% --- Title page template ---
\setbeamertemplate{title page}{%
  \vspace{1.8em}%
  \begin{beamercolorbox}[sep=6pt, center]{title}%
    \usebeamerfont{title}\usebeamercolor[fg]{title}\inserttitle\par%
  \end{beamercolorbox}%
  \vskip0.4em%
  \begin{beamercolorbox}[sep=4pt, center]{subtitle}%
    \usebeamerfont{subtitle}\usebeamercolor[fg]{subtitle}\insertsubtitle\par%
  \end{beamercolorbox}%
  \vskip0.8em%
  {\color{bathnav!40}\rule{0.5\paperwidth}{0.4pt}}\par%
  \vskip0.8em%
  \begin{beamercolorbox}[sep=3pt, center]{author}%
    \usebeamerfont{author}\insertauthor%
  \end{beamercolorbox}%
  \begin{beamercolorbox}[sep=2pt, center]{institute}%
    \usebeamerfont{institute}\insertinstitute%
  \end{beamercolorbox}%
  \vskip0.4em%
  \begin{beamercolorbox}[sep=2pt, center]{date}%
    \usebeamerfont{date}\insertdate%
  \end{beamercolorbox}%
}

% ============================================================
% SECTION TRANSITION SLIDES
% ============================================================

% Usage: \sectionslide{Section Title}{Optional subtitle or question}
\newcommand{\sectionslide}[2]{%
  \begingroup%
  \setbeamercolor{background canvas}{bg=bathnav}%
  \begin{frame}[plain, noframenumbering]%
    \vfill%
    \centering%
    {\Large\bfseries\textcolor{white}{#1}}\par%
    \vskip0.8em%
    \ifx\relax#2\relax\else%
      {\normalsize\textcolor{okabesky}{#2}}\par%
    \fi%
    \vfill%
  \end{frame}%
  \endgroup%
}

% ============================================================
% TCOLORBOX CUSTOM ENVIRONMENTS
% ============================================================

\tcbset{
  beamerbase/.style={
    enhanced,
    boxrule=0.4pt,
    arc=2pt,
    left=6pt, right=6pt, top=3pt, bottom=4pt,
    fonttitle=\bfseries\small,
    attach boxed title to top left={yshift=-0.1mm},
    boxed title style={
      arc=2pt,
      boxrule=0pt,
      left=6pt, right=6pt, top=3pt, bottom=3pt,
    },
    before upper={\setlength\parindent{0pt}\small},
  }
}

% Key point box — dark navy title bar, white title text, light blue body
\newtcolorbox{keybox}[1][Key Point]{
  beamerbase,
  colback=bathnav!8,
  colframe=bathnav,
  colbacktitle=bathnav,
  coltitle=white,
  title={#1},
}

% Paper example box — dark orange title bar, white title text, light orange body
\newtcolorbox{examplebox}[1][Example]{
  beamerbase,
  colback=okabeorange!12,
  colframe=okabeorange!80!black,
  colbacktitle=okabeorange!85!black,
  coltitle=white,
  title={#1},
}

% Warning / caveat box — vermillion title bar, white title text, light red body
\newtcolorbox{warningbox}[1][Caution]{
  beamerbase,
  colback=okabevermillion!10,
  colframe=okabevermillion,
  colbacktitle=okabevermillion,
  coltitle=white,
  title={#1},
}

% Intuition box — green title bar, white title text, light green body
\newtcolorbox{intuitionbox}[1][Intuition]{
  beamerbase,
  colback=okabegreen!10,
  colframe=okabegreen,
  colbacktitle=okabegreen,
  coltitle=white,
  title={#1},
}

% ============================================================
% MATH MACROS
% ============================================================

\newcommand{\E}[1]{\mathbb{E}\!\left[#1\right]}
\newcommand{\Econd}[2]{\mathbb{E}\!\left[#1\,\middle|\,#2\right]}
\newcommand{\Yi}[1]{Y_{i}(#1)}
\newcommand{\Y}[1]{Y(#1)}
\newcommand{\ATE}{\mathrm{ATE}}
\newcommand{\ATT}{\mathrm{ATT}}
\newcommand{\ITT}{\mathrm{ITT}}
\newcommand{\LATE}{\mathrm{LATE}}
\newcommand{\indep}{\perp\!\!\!\perp}
\newcommand{\muted}[1]{{\color{nearblack!50}#1}}
\newcommand{\highlight}[1]{{\color{bathnav}\textbf{#1}}}
\newcommand{\source}[1]{\vspace{0.3em}\par{\scriptsize\color{nearblack!60}\textit{#1}}}

% ============================================================
% BIBLIOGRAPHY STYLE
% ============================================================

\bibliographystyle{plainnat}


% ============================================================
% METADATA
% ============================================================
\title[D1B: Lab --- Settele (2022)]{Beliefs, Information, and Policy Demand}
\subtitle{Replicating \citet{Settele2022_gender}}
\author{Na Zou}
\institute{University of Bath \\ School of Management}
\date{MN52233 \quad Day 1 $\cdot$ Session B}

% ============================================================
\begin{document}
% ============================================================

% --- Title slide ---
\begin{frame}[plain, noframenumbering]
  \titlepage
\end{frame}

% ============================================================
% SESSION ROADMAP
% ============================================================

\begin{frame}[t]{Session Roadmap}
\vspace{0.4em}
\begin{center}
\footnotesize
\begin{tabular}{llr}
\toprule
\textbf{Section} & \textbf{Content} & \textbf{Approx.\ time} \\
\midrule
1. The Paper        & Research question, context, main finding     & 20 min \\
2. Research Design  & Survey flow, treatment, identification        & 25 min \\
3. Replication      & STATA hands-on: balance check + main OLS     & 55 min \\
\bottomrule
\end{tabular}
\end{center}
\vspace{0.5em}
\begin{keybox}[Lab objective]
Work through a published survey experiment from scratch:
understand the design, verify randomisation, and replicate the main causal estimate.
\end{keybox}
\vspace{0.4em}
\muted{\small Data: pre-loaded on shared drive $\cdot$
Software: STATA $\cdot$
Paper: \citet{Settele2022_gender}}
\end{frame}

% ============================================================
\sectionslide{1. The Paper}{How beliefs about the gender wage gap shape policy demand}
% ============================================================

\begin{frame}[t]{Research Question}
\begin{columns}[T]
\begin{column}{0.58\textwidth}
\begin{itemize}
  \setlength\itemsep{0.5em}
  \item The gender wage gap is a persistent feature of labour markets worldwide.
  \item Governments have policy tools available: equal pay legislation, pay transparency mandates, parental leave.
  \item \textbf{But do people support these policies — and why?}
  \item One candidate mechanism: \textbf{beliefs} about how large the gap actually is.
\end{itemize}
\vspace{0.5em}
\begin{keybox}[\citealt{Settele2022_gender}]
\textit{``How Do Beliefs about the Gender Wage Gap Affect the Demand for Public Policy?''}\\
\textit{AEJ: Economic Policy}, 14(2), 475--508.
\end{keybox}
\end{column}
\begin{column}{0.38\textwidth}
\begin{center}
\begin{tikzpicture}[scale=0.88, every node/.style={font=\small}]
  \node[draw, rounded corners, fill=okabesky!20, minimum width=2.8cm,
        minimum height=0.75cm, thick, color=okabesky!80!black] (B) at (0, 2.2) {Beliefs};
  \node[draw, rounded corners, fill=okabeblue!15, minimum width=2.8cm,
        minimum height=0.75cm, thick, color=okabeblue] (P) at (0, 0) {Policy support};
  \draw[->, thick, color=okabeblue, line width=1.2pt] (B) -- (P);
  \node[font=\scriptsize, color=nearblack!60] at (1.6, 1.1) {Causal?};
\end{tikzpicture}
\end{center}
\end{column}
\end{columns}
\end{frame}

% ---

\begin{frame}[t]{Why Beliefs? Two Channels}
\begin{columns}[T]
\begin{column}{0.48\textwidth}
\begin{examplebox}[Channel 1: Fairness concern]
If people believe women are paid significantly less for equal work, they may demand policy to correct the unfairness.
\end{examplebox}
\vspace{0.5em}
\begin{examplebox}[Channel 2: Factual updating]
People's beliefs about the gap are often wrong. Providing accurate information could shift both beliefs and policy preferences.
\end{examplebox}
\end{column}
\begin{column}{0.48\textwidth}
\begin{warningbox}[The identification problem]
Just correlating beliefs with policy support is not enough.\\[0.3em]
Those who believe the gap is large may differ in unobservable ways (political views, values) from those who think it is small.\\[0.3em]
\textbf{Selection bias} — the same problem we saw in Session~A.
\end{warningbox}
\end{column}
\end{columns}
\end{frame}

% ---

\begin{frame}[t]{What Do People Actually Believe?}
\begin{columns}[T]
\begin{column}{0.54\textwidth}
\begin{itemize}
  \setlength\itemsep{0.5em}
  \item In the US, women earn approximately \textbf{82 cents} for every dollar earned by men (unadjusted gap $\approx 18$\%).
  \item Settele surveys respondents \emph{before} any treatment and records their belief about this gap.
  \item Finding: most respondents \textbf{underestimate} the gap.
  \item Many believe women earn closer to 90--95 cents on the dollar.
  \item Implication: there is room for belief-updating.
\end{itemize}
\vspace{0.3em}
\muted{\small Heterogeneity in prior beliefs turns out to be important for \emph{who} responds to the treatment.}
\end{column}
\begin{column}{0.42\textwidth}
\begin{center}
\begin{tikzpicture}[scale=0.80, every node/.style={font=\scriptsize}]
  % Axes
  \draw[->] (0,0) -- (5.4,0) node[right] {Believed gap (\%)};
  \draw[->] (0,0) -- (0,3.4) node[above, align=center] {Share of\\respondents};
  % Schematic bars (illustrative)
  \fill[okabeblue!30] (0.3,0) rectangle (1.0,0.5);
  \fill[okabeblue!40] (1.0,0) rectangle (1.7,1.1);
  \fill[okabeblue!55] (1.7,0) rectangle (2.4,2.3);
  \fill[okabeblue!65] (2.4,0) rectangle (3.1,2.7);
  \fill[okabeblue!45] (3.1,0) rectangle (3.8,1.5);
  \fill[okabeblue!30] (3.8,0) rectangle (4.5,0.7);
  % x-axis ticks
  \node at (0.65, -0.3) {5};
  \node at (1.35, -0.3) {10};
  \node at (2.05, -0.3) {15};
  \node at (2.75, -0.3) {20};
  \node at (3.45, -0.3) {25};
  \node at (4.15, -0.3) {30+};
  % True value line
  \draw[dashed, thick, color=okabevermillion]
        (2.05,0) -- (2.05,3.1)
        node[above, color=okabevermillion] {True: 18\%};
  % Annotation
  \node[color=nearblack!55, align=center] at (3.5, 2.95)
        {\tiny Most underestimate};
\end{tikzpicture}
\end{center}
\source{Schematic based on \citet{Settele2022_gender}}
\end{column}
\end{columns}
\end{frame}

% ---

\begin{frame}[t]{The Key Insight}
\begin{itemize}
  \setlength\itemsep{0.5em}
  \item If beliefs are wrong \emph{and} beliefs drive policy preferences, then exogenously correcting beliefs should shift preferences.
  \item But we cannot just compare high-belief vs.\ low-belief respondents — they differ in many ways.
  \item Solution: \textbf{randomly assign} who receives accurate information about the gap.
\end{itemize}
\vspace{0.5em}
\begin{intuitionbox}[The experimental logic]
Randomly assign some respondents to receive factual information about the gender wage gap.\\[0.3em]
Compare their policy support to the control group who received no such information.\\[0.3em]
Any difference in outcomes is \textbf{causal} — caused by the information, not by pre-existing differences between respondents.
\end{intuitionbox}
\end{frame}

% ---

\begin{frame}[t]{Sample and Setting}
\begin{columns}[T]
\begin{column}{0.54\textwidth}
\begin{itemize}
  \setlength\itemsep{0.5em}
  \item \textbf{Platform:} Online survey (US respondents)
  \item \textbf{Sample size:} Approximately 1,500 adult respondents
  \item \textbf{Recruitment:} Broadly representative of the US adult population by gender, age, and education
  \item \textbf{Design:} Pre-registered randomised survey experiment
  \item \textbf{Ethics:} Participants consent to the survey; no deception involved
\end{itemize}
\end{column}
\begin{column}{0.42\textwidth}
\begin{keybox}[Why online surveys?]
\textbf{Advantages}
\begin{itemize}
  \setlength\itemsep{0.25em}
  \item Cost-effective at scale
  \item Rapid data collection
  \item Easy randomisation at the point of assignment
  \item Diverse national sample
\end{itemize}
\vspace{0.3em}
\textbf{Limitation}
\begin{itemize}
  \setlength\itemsep{0.25em}
  \item Hypothetical policy preferences may differ from real voting behaviour
\end{itemize}
\end{keybox}
\end{column}
\end{columns}
\end{frame}

% ---

\begin{frame}[t]{What Settele Finds}
\begin{itemize}
  \setlength\itemsep{0.5em}
  \item \textbf{Belief updating:} The information treatment significantly shifts respondents' beliefs toward the true gender wage gap.
  \item \textbf{Policy demand:} Respondents who receive the information show \emph{higher} support for equal-pay policies.
  \item \textbf{Mechanism:} The effect runs through beliefs — not just general political priming or awareness.
  \item \textbf{Heterogeneity:} Effects are concentrated among those who initially \emph{underestimated} the gap.
\end{itemize}
\vspace{0.5em}
\begin{examplebox}[The main result we will replicate]
Providing information about the gender wage gap \textbf{causally increases} support for equal-pay policies
by a statistically and economically significant amount.
\end{examplebox}
\end{frame}

% ---

\begin{frame}[t]{Connecting to Session A}
\begin{columns}[T]
\begin{column}{0.55\textwidth}
\begin{itemize}
  \setlength\itemsep{0.5em}
  \item This is a \textbf{survey experiment} — treatment is an information provision, not a physical intervention.
  \item Random assignment ensures $D_i \indep (Y_i(0), Y_i(1))$.
  \item With full compliance (no opt-out from reading), $\ITT = \ATE$.
  \item The SUTVA assumption holds: one respondent's treatment does not affect another respondent's outcome.
\end{itemize}
\end{column}
\begin{column}{0.42\textwidth}
\begin{keybox}[From D1A to D1B]
\small
\begin{tabular}{ll}
\toprule
D1A concept & D1B instance \\
\midrule
Treatment $D_i$ & Information provision \\
Outcome $Y_i$ & Policy support index \\
$\ATE$ & Avg.\ effect on support \\
Randomisation & Survey random assignment \\
\bottomrule
\end{tabular}
\end{keybox}
\end{column}
\end{columns}
\end{frame}

% ============================================================
\sectionslide{2. Research Design}{Survey flow, treatment, and identification}
% ============================================================

\begin{frame}[t]{Survey Flow: Three Stages}
\vspace{0.3em}
\begin{center}
\begin{tikzpicture}[scale=0.90, every node/.style={font=\small}]
  % Stage 1
  \node[draw, rounded corners, fill=okabesky!20, minimum width=3.2cm, minimum height=1.0cm,
        thick, color=okabesky!80!black, align=center] (S1) at (0, 0)
        {\textbf{Stage 1}\\Prior belief elicitation};
  % Stage 2
  \node[draw, rounded corners, fill=okabeorange!20, minimum width=3.2cm, minimum height=1.0cm,
        thick, color=okabeorange!90!black, align=center] (S2) at (5, 0)
        {\textbf{Stage 2}\\Randomised treatment};
  % Stage 3
  \node[draw, rounded corners, fill=okabeblue!15, minimum width=3.2cm, minimum height=1.0cm,
        thick, color=okabeblue, align=center] (S3) at (10, 0)
        {\textbf{Stage 3}\\Policy outcome};
  % Arrows
  \draw[->, thick, color=nearblack!60] (S1) -- (S2);
  \draw[->, thick, color=nearblack!60] (S2) -- (S3);
  % Sub-labels
  \node[font=\scriptsize, color=nearblack!60, align=center] at (0, -0.95)
        {``What \% do women earn\\per \$1 men earn?''};
  \node[font=\scriptsize, color=okabeorange!90!black, align=center] at (5, -0.95)
        {Treatment: actual gap stats\\Control: unrelated filler};
  \node[font=\scriptsize, color=nearblack!60, align=center] at (10, -0.95)
        {Support for equal pay,\\pay transparency, leave};
  % Randomisation annotation
  \node[font=\scriptsize, color=bathnav, align=center] at (5, 0.95)
        {$\uparrow$ Random assignment};
\end{tikzpicture}
\end{center}
\vspace{0.4em}
\begin{keybox}[Critical design feature]
Prior beliefs are elicited \emph{before} randomisation.
This ensures the belief measure is not contaminated by the treatment.
\end{keybox}
\end{frame}

% ---

\begin{frame}[t]{The Information Treatment}
\begin{columns}[T]
\begin{column}{0.54\textwidth}
\textbf{Treatment group receives:}
\begin{itemize}
  \setlength\itemsep{0.4em}
  \item The official BLS/Census estimate of the unadjusted gender wage gap ($\approx 18$\%)
  \item Framing: ``Women earn approximately 82 cents for every dollar earned by men''
  \item Brief factual context; no advocacy or political framing
\end{itemize}
\vspace{0.4em}
\textbf{Control group receives:}
\begin{itemize}
  \setlength\itemsep{0.4em}
  \item Unrelated filler content on a different topic
  \item Matched in length and format to avoid differential survey fatigue
\end{itemize}
\end{column}
\begin{column}{0.42\textwidth}
\begin{warningbox}[Demand effects]
Does the treatment inadvertently signal the ``right'' answer to give?\\[0.3em]
Settele addresses this with:\\[0.2em]
$\cdot$ Neutral, factual framing\\
$\cdot$ Outcome questions on a separate screen\\
$\cdot$ Robustness checks using political priming placebo
\end{warningbox}
\end{column}
\end{columns}
\end{frame}

% ---

\begin{frame}[t]{Outcome Variables}
\begin{columns}[T]
\begin{column}{0.54\textwidth}
\begin{itemize}
  \setlength\itemsep{0.4em}
  \item Respondents rate their support for several policies on a Likert scale.
  \item Settele aggregates these into a \textbf{policy support index} (standardised):
  \begin{itemize}
    \setlength\itemsep{0.25em}
    \item Equal pay legislation
    \item Pay transparency requirements
    \item Government-mandated parental leave
    \item Affirmative action in hiring
  \end{itemize}
  \item Higher score = stronger support for gender-targeted policies.
\end{itemize}
\end{column}
\begin{column}{0.42\textwidth}
\begin{keybox}[Why an index?]
Individual survey items are noisy.\\[0.2em]
Averaging across related items:
\begin{itemize}
  \setlength\itemsep{0.25em}
  \item Reduces measurement error
  \item Increases statistical power
  \item Guards against multiple testing concerns
\end{itemize}
\vspace{0.3em}
\muted{\small Standard practice in survey experiments\\
\citep{Stantcheva2023_surveys}}
\end{keybox}
\end{column}
\end{columns}
\end{frame}

% ---

\begin{frame}[t]{Identification: Why Randomisation Works}
\begin{itemize}
  \setlength\itemsep{0.5em}
  \item Random assignment creates two groups that are, \textbf{on average}, identical in all pre-treatment characteristics.
  \item Any post-treatment difference in policy support is therefore attributable to the treatment alone.
\end{itemize}
\vspace{0.3em}
\begin{keybox}[Identification assumption]
\[
  D_i \indep \bigl(Y_i(0),\, Y_i(1)\bigr)
\]
Treatment assignment is independent of potential outcomes.
This holds \textbf{by design} under random assignment — not as an untestable assumption.
\end{keybox}
\vspace{0.3em}
\muted{\small We verify this empirically in Section~3 via a balance table (Step~4).}
\end{frame}

% ---

\begin{frame}[t]{Potential Outcomes in This Setting}
\begin{columns}[T]
\begin{column}{0.55\textwidth}
For respondent $i$:
\begin{align*}
  Y_i(1) &= \text{policy support \emph{if} given information} \\
  Y_i(0) &= \text{policy support \emph{if not} given information}
\end{align*}
\begin{itemize}
  \setlength\itemsep{0.4em}
  \item We observe only $Y_i = D_i \cdot Y_i(1) + (1-D_i) \cdot Y_i(0)$.
  \item \textbf{Estimand:} Average Treatment Effect
    \[
      \ATE = \E{Y_i(1) - Y_i(0)}
    \]
  \item Randomisation allows us to estimate $\ATE$ with a simple difference in means.
\end{itemize}
\end{column}
\begin{column}{0.41\textwidth}
\begin{intuitionbox}[What we compute]
\[
  \widehat{\ATE} = \bar{Y}_{\mathrm{treat}} - \bar{Y}_{\mathrm{control}}
\]
Under randomisation:
\[
\Econd{Y_i(0)}{D_i=1} = \Econd{Y_i(0)}{D_i=0}
\]
so the control group is a valid counterfactual for the treatment group.
\end{intuitionbox}
\end{column}
\end{columns}
\end{frame}

% ---

\begin{frame}[t]{Why Not Just Correlate Beliefs with Policy Support?}
\begin{columns}[T]
\begin{column}{0.55\textwidth}
\textbf{Observational approach:}
\begin{itemize}
  \setlength\itemsep{0.4em}
  \item Compare policy support of respondents who \emph{already believe} the gap is large vs.\ those who think it is small.
  \item Problem: high-belief respondents may be more left-leaning, more educated, more likely to identify as feminist.
  \item These confounders — not beliefs per se — could drive the policy support difference.
\end{itemize}
\end{column}
\begin{column}{0.41\textwidth}
\begin{warningbox}[Selection problem]
\begin{center}
\begin{tikzpicture}[scale=0.78, every node/.style={font=\scriptsize}]
  \node[draw, rounded corners, fill=okabepink!20, minimum width=2.1cm,
        minimum height=0.65cm, color=okabepink!80!black] (C) at (0, 2.1) {Politics, education};
  \node[draw, rounded corners, fill=okabesky!20, minimum width=1.9cm,
        minimum height=0.65cm] (B) at (-1.5, 0) {High beliefs};
  \node[draw, rounded corners, fill=okabeblue!15, minimum width=1.9cm,
        minimum height=0.65cm] (P) at (1.5, 0) {Policy support};
  \draw[->, thick, color=okabepink!70!black] (C) -- (B);
  \draw[->, thick, color=okabepink!70!black] (C) -- (P);
  \draw[<->, dashed, thick, color=nearblack!40]
        (B) -- node[below, color=nearblack!50]{spurious?}(P);
\end{tikzpicture}
\end{center}
\end{warningbox}
\end{column}
\end{columns}
\end{frame}

% ---

\begin{frame}[t]{The Estimator: OLS with Controls}
\begin{itemize}
  \setlength\itemsep{0.5em}
  \item The causal estimator is a simple OLS regression of policy support on the treatment indicator.
\end{itemize}
\vspace{0.3em}
\begin{keybox}[Regression specification]
\[
  Y_i = \alpha + \beta \, D_i + \mathbf{X}_i' \gamma + \varepsilon_i
\]
\vspace{-0.5em}
\begin{tabular}{ll}
  $Y_i$ & policy support index \\
  $D_i$ & treatment indicator ($1 =$ information received) \\
  $\mathbf{X}_i$ & pre-treatment controls (gender, age, education, income, party ID) \\
  $\beta$ & \textbf{ATE estimate} --- the coefficient of interest \\
\end{tabular}
\end{keybox}
\vspace{0.3em}
\muted{\small Controls improve precision but do not change the interpretation of $\hat\beta$ under randomisation.}
\end{frame}

% ---

\begin{frame}[t]{Threats to Validity}
\begin{columns}[T]
\begin{column}{0.48\textwidth}
\begin{warningbox}[Internal validity threats]
\begin{itemize}
  \setlength\itemsep{0.3em}
  \item \textbf{Demand effects:} respondents guess what the researcher wants to hear
  \item \textbf{Differential attrition:} treatment group drops out at different rates
  \item \textbf{Non-compliance:} respondents skip or skim the treatment content
\end{itemize}
\end{warningbox}
\end{column}
\begin{column}{0.48\textwidth}
\begin{warningbox}[External validity threats]
\begin{itemize}
  \setlength\itemsep{0.3em}
  \item \textbf{Online sample:} may differ from the general population
  \item \textbf{Hypothetical responses:} stated vs.\ revealed preferences
  \item \textbf{One-time exposure:} real-world information arrives repeatedly from multiple sources
\end{itemize}
\end{warningbox}
\end{column}
\end{columns}
\vspace{0.5em}
\muted{\small Settele tests robustness to demand effects and attrition. These checks are standard practice in survey experiments.}
\end{frame}

% ============================================================
\sectionslide{3. Hands-On Replication}{STATA: balance check + main treatment effect}
% ============================================================

\begin{frame}[t]{Data Overview}
\begin{columns}[T]
\begin{column}{0.56\textwidth}
\textbf{Dataset:} \texttt{settele2022.dta} (pre-loaded on shared drive)
\vspace{0.5em}

\begin{footnotesize}
\begin{tabular}{lll}
\toprule
\textbf{Variable} & \textbf{Type} & \textbf{Description} \\
\midrule
\texttt{treat}         & Binary  & 1 = information treatment \\
\texttt{policy\_index} & Cont.   & Policy support (standardised) \\
\texttt{belief\_pre}   & Cont.   & Prior belief: wage gap (\%) \\
\texttt{belief\_post}  & Cont.   & Post-treatment belief (\%) \\
\texttt{female}        & Binary  & 1 = female respondent \\
\texttt{age}           & Cont.   & Age in years \\
\texttt{educ}          & Ordinal & Education level (1--5) \\
\texttt{income}        & Ordinal & Household income category (1--7) \\
\texttt{pid}           & Nominal & Party ID (1=Dem, 2=Ind, 3=Rep) \\
\bottomrule
\end{tabular}
\end{footnotesize}
\end{column}
\begin{column}{0.40\textwidth}
\begin{keybox}[Codebook check]
\small
Variable names above are illustrative.\\[0.3em]
\textbf{Always verify} names against the dataset:
\vspace{0.3em}
\begin{itemize}
  \setlength\itemsep{0.25em}
  \item \texttt{describe} — all variables and labels
  \item \texttt{codebook} — ranges, types, labels
  \item \texttt{README.txt} — data documentation
\end{itemize}
\vspace{0.3em}
\muted{\small Replication packages from AEJ always include a README.}
\end{keybox}
\end{column}
\end{columns}
\end{frame}

% ---

\begin{frame}[fragile, t]{Step 0: Setup}
Set up your STATA environment before doing anything else.
\vspace{0.3em}
\begin{examplebox}[What this achieves]
Clears memory, points STATA to the shared drive, loads the data, and opens a log file for reproducibility.
\end{examplebox}
\vspace{0.3em}
\begin{footnotesize}
\begin{verbatim}
*--- Step 0: Setup ---*
clear all
set more off

* Set working directory (adjust path to your shared drive)
cd "Z:/MN52233/Lab_D1B"

* Open log file
log using "settele_replication.log", replace text

* Load data
use "settele2022.dta", clear

* Quick sanity check: how many observations?
count
\end{verbatim}
\end{footnotesize}
\end{frame}

% ---

\begin{frame}[fragile, t]{Step 1: Explore the Data}
Always look at the data before running regressions.
\vspace{0.3em}
\begin{examplebox}[What this achieves]
Understand structure, check for missing values, and confirm the treatment split looks as expected.
\end{examplebox}
\vspace{0.3em}
\begin{footnotesize}
\begin{verbatim}
*--- Step 1: Explore ---*

* Variable overview with labels
describe

* Summary statistics for key variables
summarize treat policy_index belief_pre belief_post ///
          female age educ income

* How many respondents in each condition?
tab treat

* Flag any missing values
misstable summarize treat policy_index belief_pre female age educ income
\end{verbatim}
\end{footnotesize}
\end{frame}

% ---

\begin{frame}[fragile, t]{Step 2: Visualise Prior Beliefs}
Plot the distribution of respondents' beliefs about the gender wage gap \emph{before} any treatment.
Where does the true value fall relative to the modal guess?
\vspace{0.4em}
\begin{footnotesize}
\begin{verbatim}
*--- Step 2: Visualise prior beliefs ---*
histogram belief_pre, ///
    bin(20) ///
    color("0 114 178") lcolor(white) ///
    xline(18, lcolor("213 94 0") lwidth(medthick) lpattern(dash)) ///
    xtitle("Believed gender wage gap (%)") ///
    ytitle("Density") ///
    title("Prior beliefs about the gender wage gap", size(medium)) ///
    note("Dashed line = actual gap (~18%). Source: Settele (2022)", size(small)) ///
    graphregion(color(white)) bgcolor(white)

graph export "fig_belief_prior.png", replace width(1200)
\end{verbatim}
\end{footnotesize}
\vspace{0.2em}
\muted{\small \textbf{Look for:} Is the distribution concentrated below 18\%? That would confirm systematic underestimation.}
\end{frame}

% ---

\begin{frame}[fragile, t]{Step 3: Does the Treatment Update Beliefs?}
Before looking at the policy outcome, check whether the treatment actually changes what people believe.
\begin{examplebox}[This is not the main result]
This is a check on the mechanism: does the information \emph{land}?
\end{examplebox}
\vspace{0.3em}
\begin{footnotesize}
\begin{verbatim}
*--- Step 3: Belief updating ---*

* Compare post-treatment beliefs between groups
ttest belief_post, by(treat)

* Regression version (equivalent, more flexible)
reg belief_post treat, robust

* How much did beliefs shift?
gen belief_change = belief_post - belief_pre
ttest belief_change, by(treat)
\end{verbatim}
\end{footnotesize}
\vspace{0.2em}
\muted{\small \textbf{Expect:} The treatment group shows a significantly higher \texttt{belief\_post} and larger \texttt{belief\_change}.}
\end{frame}

% ---

\begin{frame}[t]{Step 4: Randomisation Check --- Theory}
\begin{itemize}
  \setlength\itemsep{0.5em}
  \item Random assignment should produce treatment and control groups that are \textbf{statistically similar} on all pre-treatment characteristics.
  \item If we find large, systematic differences in baseline variables, something may have gone wrong with the randomisation.
\end{itemize}
\vspace{0.4em}
\begin{keybox}[What to check]
For each baseline variable $X_k$, estimate:
\[
  X_k = \alpha_k + \delta_k D_i + u_i
\]
If $\hat\delta_k \approx 0$ and $p > 0.10$ for all $k$: \textbf{balance confirmed.}\\[0.2em]
One or two marginally significant results are expected by chance alone.
\end{keybox}
\vspace{0.3em}
\muted{\small This is the ``balance table'' or ``Table~1'' found in almost every experimental paper.}
\end{frame}

% ---

\begin{frame}[fragile, t]{Step 5: Balance Table}
\begin{footnotesize}
\begin{verbatim}
*--- Step 5: Randomisation check ---*

* Regress each baseline covariate on treatment
foreach var in female age educ income {
    reg `var' treat, robust
    estimates store bal_`var'
}

* Display results side by side
esttab bal_female bal_age bal_educ bal_income, ///
    cells(b(fmt(3)) se(fmt(3) par)) ///
    stats(N r2, fmt(0 3) labels("N" "R-squared")) ///
    mtitles("Female" "Age" "Education" "Income") ///
    title("Table: Baseline balance across treatment and control") ///
    nostar

* Alternative: simple t-tests
foreach var in female age educ income {
    ttest `var', by(treat)
}
\end{verbatim}
\end{footnotesize}
\vspace{0.2em}
\muted{\small \textbf{Expect:} all coefficients close to zero, $p > 0.10$, low $R^2$.}
\end{frame}

% ---

\begin{frame}[fragile, t]{Step 6: Main Regression}
Now estimate the treatment effect on the policy support index.
\vspace{0.2em}
\begin{keybox}[Specification]
$Y_i = \alpha + \beta\,D_i + \gamma_1\,\texttt{female}_i + \gamma_2\,\texttt{age}_i + \gamma_3\,\texttt{educ}_i + \gamma_4\,\texttt{income}_i + \boldsymbol{\delta}'\mathbf{1}[\texttt{pid}_i] + \varepsilon_i$
\end{keybox}
\vspace{0.3em}
\begin{footnotesize}
\begin{verbatim}
*--- Step 6: Main treatment effect ---*

* Model 1: unadjusted (treatment only)
reg policy_index treat, robust
estimates store m1

* Model 2: with pre-treatment controls
reg policy_index treat female age educ income i.pid, robust
estimates store m2

* Export results table
esttab m1 m2 using "table_main.rtf", ///
    cells(b(fmt(3) star) se(fmt(3) par)) ///
    stats(N r2, fmt(0 3) labels("N" "R-squared")) ///
    mtitles("No controls" "With controls") ///
    title("Table 1: Treatment effect on policy support index") ///
    replace
\end{verbatim}
\end{footnotesize}
\end{frame}

% ---

\begin{frame}[t]{Step 7: Reading the Output}
\begin{columns}[T]
\begin{column}{0.54\textwidth}
The \textbf{coefficient on \texttt{treat}} ($\hat\beta$) tells us:
\begin{itemize}
  \setlength\itemsep{0.4em}
  \item The average difference in policy support index between treatment and control groups.
  \item Since the index is \emph{standardised}, $\hat\beta$ is measured in \textbf{standard deviation units}.
  \item A $\hat\beta = 0.10$ means the treatment shifts policy support by 0.10~SD.
\end{itemize}
\vspace{0.4em}
\begin{keybox}[Questions to ask]
Is $\hat\beta$ positive?\\
Is it statistically significant?\\
Does adding controls change $\hat\beta$ much?\\
\muted{\small If not: randomisation worked as intended.}
\end{keybox}
\end{column}
\begin{column}{0.42\textwidth}
\begin{examplebox}[Illustrative output]
\small
\begin{tabular}{lcc}
\toprule
 & (1) & (2) \\
\midrule
\texttt{treat} & $0.11^{**}$ & $0.10^{**}$ \\
               & $(0.05)$    & $(0.05)$    \\
Controls       & No          & Yes         \\
$N$            & 1,512       & 1,512       \\
$R^2$          & 0.006       & 0.182       \\
\bottomrule
\multicolumn{3}{l}{\tiny $^{**}p<0.05$. Robust SEs in parentheses.}
\end{tabular}
\end{examplebox}
\vspace{0.3em}
\muted{\small Your actual results will differ. The key feature to look for is the sign and significance of column~(1).}
\end{column}
\end{columns}
\end{frame}

% ---

\begin{frame}[fragile, t]{Step 8: Visualise the Result}
A bar chart of group means is often the clearest way to communicate the treatment effect.
\vspace{0.4em}
\begin{footnotesize}
\begin{verbatim}
*--- Step 8: Visualise the treatment effect ---*

* Bar chart: mean policy support by treatment group
graph bar (mean) policy_index, over(treat, ///
        relabel(1 "Control" 2 "Treatment")) ///
    bar(1, color("86 180 233")) ///
    bar(2, color("0 114 178")) ///
    ytitle("Policy support index (SD units)") ///
    title("Figure: Mean policy support by treatment group", size(medium)) ///
    note("Error bars: 95% CIs. Source: Settele (2022)", size(small)) ///
    blabel(bar, format(%5.3f) size(small)) ///
    graphregion(color(white)) bgcolor(white)

graph export "fig_main_result.png", replace width(1200)
\end{verbatim}
\end{footnotesize}
\end{frame}

% ---

\begin{frame}[t]{Interpretation: Back to Potential Outcomes}
\begin{itemize}
  \setlength\itemsep{0.5em}
  \item Your estimated $\hat\beta$ is the sample analogue of:
    \[
      \ATE = \E{Y_i(1)} - \E{Y_i(0)}
    \]
  \item Under randomisation, $\bar{Y}_{\mathrm{treat}}$ estimates $\E{Y_i(1)}$ and $\bar{Y}_{\mathrm{control}}$ estimates $\E{Y_i(0)}$.
  \item Adding controls $\mathbf{X}_i$ reduces residual variance and tightens standard errors.
    It does \emph{not} change the target estimand $\ATE$ under randomisation.
\end{itemize}
\vspace{0.3em}
\begin{intuitionbox}[The punchline]
Information about the gender wage gap \textbf{causally increases} support for equal-pay policies.\\[0.2em]
This is not a correlation — it is identified by the \textbf{randomisation} of who received the information.
\end{intuitionbox}
\end{frame}

% ---

\begin{frame}[t]{Discussion Questions}
\vspace{0.3em}
\begin{columns}[T]
\begin{column}{0.48\textwidth}
\begin{examplebox}[Q1: External validity]
Our sample is US online respondents.\\[0.3em]
Would you expect the same result in a country with a different gender wage gap, a different policy environment, or different baseline beliefs?\\[0.2em]
What evidence would help you answer this?
\end{examplebox}
\vspace{0.4em}
\begin{examplebox}[Q2: Mechanism]
We know the treatment shifts beliefs and policy support. But \emph{how} does a belief change translate into a policy preference?\\[0.2em]
Fairness concerns? Economic reasoning? Social signalling?
\end{examplebox}
\end{column}
\begin{column}{0.48\textwidth}
\begin{examplebox}[Q3: Policy implications]
Suppose a government ran a public information campaign about the gender wage gap.\\[0.3em]
Based on Settele's findings, what would you predict?\\[0.2em]
What additional evidence would you want before recommending such a campaign?
\end{examplebox}
\end{column}
\end{columns}
\end{frame}

% ---

\begin{frame}[t]{Wrap-Up}
\begin{columns}[T]
\begin{column}{0.54\textwidth}
\textbf{What we covered today:}
\begin{itemize}
  \setlength\itemsep{0.4em}
  \item A real survey experiment (\textit{AEJ: EP}, 2022)
  \item How information provision creates causal identification
  \item Randomisation check: balance table in STATA
  \item Main OLS treatment effect with and without controls
  \item Visualising and interpreting experimental results in STATA
\end{itemize}
\vspace{0.4em}
\begin{keybox}[Key takeaway]
In a well-designed survey experiment, a simple OLS of outcome on treatment recovers the ATE.
The hard work is in the \emph{design}, not the regression.
\end{keybox}
\end{column}
\begin{column}{0.42\textwidth}
\textbf{Coming up on Day 2:}
\begin{itemize}
  \setlength\itemsep{0.4em}
  \item \textbf{Session A:} Regression Discontinuity and Difference-in-Differences
  \item \textbf{Session B:} Replication of \citet{Braghieri2022_socialmedia}
  \begin{itemize}
    \item Facebook rollout as a natural experiment
    \item Staggered DiD design
  \end{itemize}
\end{itemize}
\vspace{0.5em}
\muted{\small Preparation: read the empirical strategy section of \citealt{Braghieri2022_socialmedia}}
\end{column}
\end{columns}
\end{frame}

% ============================================================
% BIBLIOGRAPHY
% ============================================================

\begin{frame}[allowframebreaks]{References}
  \bibliography{../Bibliography_base}
\end{frame}

\end{document}
